% Самостійний документ. Компіляція: xelatex (двічі).
% Лекція 5: 60 хв викладу, запитання поза цим часом.
\documentclass[aspectratio=169,11pt]{beamer}
\usepackage{fontspec,polyglossia}
\setdefaultlanguage{ukrainian}
\setsansfont{DejaVu Sans}
\setmainfont{DejaVu Serif}
\setmonofont{DejaVu Sans Mono}
\newfontfamily\cyrillicfont{DejaVu Serif}
\newfontfamily\cyrillicfontsf{DejaVu Sans}
\newfontfamily\cyrillicfonttt{DejaVu Sans Mono}
\usepackage{amsmath,amssymb,booktabs,array}
\usefonttheme{professionalfonts}
\definecolor{ink}{HTML}{183341}
\definecolor{accent}{HTML}{087C83}
\setbeamercolor{normal text}{fg=ink,bg=white}
\setbeamercolor{structure}{fg=accent}
\setbeamercolor{frametitle}{fg=ink}
\setbeamertemplate{navigation symbols}{}
\setbeamertemplate{footline}{\hfill\insertframenumber\hspace{5mm}\vspace{3mm}}
\setbeamersize{text margin left=8mm,text margin right=8mm}
\setbeamerfont{frametitle}{size=\large,series=\bfseries}
\hypersetup{colorlinks=true,urlcolor=accent,linkcolor=ink}
\newcommand{\unit}[1]{\,\text{#1}}
\newcommand{\src}[2]{\par\vfill{\tiny\href{#1}{#2}}}
\newcommand{\E}{\operatorname{E}}
\newcommand{\Var}{\operatorname{Var}}
\newcommand{\Cov}{\operatorname{Cov}}
\title{Підсумовування, нормування\newline і похибки ЗВТ}
\subtitle{Лекція 5. Математичний конспект і практичні завдання}
\author{Метрологія, технологічні вимірювання та прилади}
\date{}
\begin{document}
\begin{frame}\titlepage\end{frame}

\begin{frame}{Позначення і зв’язок із лекцією 4}
\begin{itemize}
\item $b_i$ — знакове зміщення або його явно позначена оцінка; $B_i\ge0$ — межа модуля невідомого внеску.
\item $\varepsilon_i$ — випадковий внесок із нульовим середнім; $\sigma_i$ — його стандартне відхилення.
\item $c_i$ переносить внесок до одиниці кінцевого результату.
\item Паспортна межа, стандартне відхилення та стандартна невизначеність не є взаємозамінними поняттями.
\item Серія повторень може оцінювати розсіювання, але сама не охоплює всіх інструментальних впливів.
\end{itemize}
Усі числові моделі навчальні, крім явно названих параметрів Fluke 117 та WIKA S-20. Питання і пошук відповідей не входять у 60 хв викладу.
\end{frame}

\section{Як складаються похибки}
\begin{frame}{Визначення: приведення внесків до результату}
Для лінійної моделі внесок переносимо коефіцієнтом $c_i$:
\[
e_i=c_i e_{i,\mathrm{in}},\qquad b_i=c_i b_{i,\mathrm{in}},\qquad
B_i=|c_i|B_{i,\mathrm{in}},\quad \sigma_i=|c_i|\sigma_{i,\mathrm{in}}.
\]
Навчальний канал: $p\in[0;100]\unit{кПа}$, $V\in[0;10]\unit{В}$,
\[
V=Sp,\quad S=0{,}1\unit{В/кПа},\qquad
\Delta p=\frac{\Delta V}{S},\quad c=10\unit{кПа/В}.
\]
Тому $0{,}03\unit{В}$ відповідає $0{,}30\unit{кПа}$.
\medskip

Спільна одиниця є необхідною умовою додавання. Для нелінійної моделі лінійне перенесення може бути лише локальним наближенням.
\end{frame}

\begin{frame}{Визначення: алгебраїчне й граничне складання}
Якщо знакові внески відомі:
\[
b_\Sigma=\sum_i b_i,\qquad C=-b_\Sigma.
\]
$C$ — поправка; якщо $b_i$ лише оцінено, залишкову невизначеність поправки потрібно описати окремо.
\medskip

Якщо відомі тільки $|e_i|\le B_i$:
\[
\left|\sum_i e_i\right|\le\sum_i|e_i|\le\sum_iB_i=B_\Sigma.
\]
\textbf{Приклад:} $+0{,}4-0{,}3=+0{,}1\unit{кПа}$ для відомих знаків; при відомих лише межах $0{,}4$ і $0{,}3\unit{кПа}$ отримуємо $B_\Sigma=0{,}7\unit{кПа}$.
\medskip

Гранична нерівність не потребує незалежності чи припущення про розподіл.
\end{frame}

\begin{frame}{Визначення: дисперсія випадкової суми}
За існування скінченних других моментів:
\[
\varepsilon_\Sigma=\sum_i\varepsilon_i,\qquad
\Var(\varepsilon_\Sigma)=\sum_i\sigma_i^2+2\sum_{i<j}\Cov(\varepsilon_i,\varepsilon_j).
\]
Для некорельованих внесків:
\[
\sigma_\Sigma=\sqrt{\sum_i\sigma_i^2}\quad\text{(RSS)}.
\]
Незалежність є достатньою умовою нульових коваріацій. Нульова коваріація загалом не доводить незалежності.
\medskip

Для двох внесків із $\sigma_1,\sigma_2>0$:
\[
\rho=\frac{\Cov(\varepsilon_1,\varepsilon_2)}{\sigma_1\sigma_2},\quad
\sigma_\Sigma^2=\sigma_1^2+\sigma_2^2+2\rho\sigma_1\sigma_2,\quad -1\le\rho\le1.
\]
\src{https://www.nist.gov/pml/nist-technical-note-1297/nist-tn-1297-appendix-law-propagation-uncertainty}{NIST TN 1297, Appendix A: коефіцієнти впливу та коваріації в моделі невизначеності}
\end{frame}

\begin{frame}{Приклад: залежність змінює підсумок}
Для $\sigma_1=3\unit{мВ}$, $\sigma_2=4\unit{мВ}$:
\[
\sigma_\Sigma=\sqrt{25+24\rho}\unit{мВ}.
\]
\begin{center}
\begin{tabular}{ccc}
\toprule
$\rho$ & $\sigma_\Sigma$, мВ & Зміст\\\midrule
$0$ & $5$ & Некорельовані внески\\
$+1$ & $7$ & Повна додатна кореляція\\
$-1$ & $1$ & Повна від’ємна кореляція\\\bottomrule
\end{tabular}
\end{center}
Спільні живлення або температура можуть створити залежність. Різні назви приладів не доводять незалежності.
\medskip

RSS стандартних відхилень є стандартним відхиленням суми, а не її гарантованим максимумом.
\end{frame}

\begin{frame}{Визначення: повна похибка та покриття}
Нехай $\Delta=b+\varepsilon_\Sigma$, $|b|\le B$, $\E[\varepsilon_\Sigma]=0$.
\[
\{|\varepsilon_\Sigma|\le k\sigma_\Sigma\}
\subseteq\{|\Delta|\le B+k\sigma_\Sigma\}.
\]
Тому
\[
P(|\Delta|\le B+k\sigma_\Sigma)\ge P(|\varepsilon_\Sigma|\le k\sigma_\Sigma).
\]
Якщо $\varepsilon_\Sigma\sim\mathcal N(0,\sigma_\Sigma^2)$ і $\sigma_\Sigma$ відома, то $k\approx1{,}96$ дає покриття близько 95\% для випадкової складової та не менше цього для розширеної межі.
\medskip

Це модельна оцінка для результату, а не безумовна гарантія чи автоматично CI середнього. Оцінювання $\sigma$ за малою вибіркою потребує іншого обґрунтування множника.
\end{frame}

\begin{frame}{Перехід від межі до стандартної характеристики}
Сам запис $|e|\le B$ не визначає розподілу $e$.
\medskip

Лише за обґрунтованої моделі $e\sim U(-B;B)$:
\[
\E[e]=0,\qquad \sigma_e=\frac{B}{\sqrt3}.
\]
\begin{itemize}
\item Рівномірний розподіл потрібно обґрунтувати, а не вивести з символу «±».
\item У бюджеті невизначеності аналогічне перетворення може описувати стан знання про величину; це не твердження про фізичний шум приладу.
\item Нормування встановлює форму характеристики, її значення та умови; читач має розуміти, що саме нормовано.
\end{itemize}
У програмних задачах межі та СКВ зберігаємо окремо.
\end{frame}

\section{Нормовані метрологічні характеристики}
\begin{frame}{Визначення: норма та область чинності}
Нормування встановлює вимоги до характеристик ЗВТ, які впливають на результат і похибку.
\[
y=f_{\mathrm{nom}}(x),\qquad |e|\le B(x)\quad\text{за умов }\Omega.
\]
\begin{itemize}
\item $f_{\mathrm{nom}}$ задає номінальний зв’язок входу та виходу.
\item Діапазон і чутливість визначають робочі значення й масштаб перенесення внесків.
\item Варіація, межі похибок та впливні характеристики доповнюють модель відповідно до задачі.
\item $\Omega$ включає режим, діапазон, температуру, живлення, навантаження та інші заявлені умови.
\end{itemize}
Поза $\Omega$ формула може обчислюватися, але її паспортне обґрунтування відсутнє.
\src{https://jcgm.bipm.org/vim/en/4.26.html}{VIM 4.26: maximum permissible measurement error}
\end{frame}

\begin{frame}{Визначення: поширений запис паспортної межі}
\[
B(x)=a|x|+bX_{\mathrm{span}}+d q.
\]
\begin{itemize}
\item $a,b\ge0$ — безрозмірні частки: $0{,}5\%=0{,}005$.
\item $X_{\mathrm{span}}$ — ширина заданого діапазону; $q>0$ — крок молодшого розряду в одиницях величини.
\item $d\ge0$ — кількість counts; кожний доданок має одиницю $x$.
\item Використовуємо лише ті доданки, які задає конкретний паспорт.
\end{itemize}
Для $x\ne0$ відносна межа за цією моделлю:
\[
D(x)=\frac{B(x)}{|x|}\cdot100\%.
\]
При $x=0$ ця відносна форма не визначена, хоча $B(0)$ може бути додатною.
\end{frame}

\begin{frame}{Реальний приклад: Fluke 117, DC V}
Специфікація виробника: $\pm(0{,}5\%\ \text{показу}+2\ \text{counts})$.
\[
B(V)=0{,}005|V|+2q.
\]
Для навчального показу $V=5\unit{В}$:
\begin{center}
\begin{tabular}{ccc}
\toprule
Діапазон & Крок $q$ & Межа $B$\\\midrule
$6{,}000\unit{В}$ & $0{,}001\unit{В}$ & $0{,}027\unit{В}$\\
$60{,}00\unit{В}$ & $0{,}01\unit{В}$ & $0{,}045\unit{В}$\\\bottomrule
\end{tabular}
\end{center}
Це розрахунок за специфікацією за умов її застосовності, а не калібрування конкретного екземпляра. Кількість цифр дисплея не визначає межу похибки сама по собі.
\src{https://www.fluke.com/en-us/product/electrical-testing/digital-multimeters/fluke-117}{Fluke 117: DC volts, range/resolution та accuracy specification}
\end{frame}

\section{Похибки засобів вимірювань}
\begin{frame}{Визначення: класифікаційні ознаки}
\small
\begin{tabular}{>{\raggedright\arraybackslash}p{0.32\textwidth}>{\raggedright\arraybackslash}p{0.59\textwidth}}
\toprule
Поведінка повторень & Систематична або випадкова складова.\\[2mm]
Зміна входу в часі & Статична або динамічна похибка.\\[2mm]
Умови застосування & Основна за нормальних умов; додаткова через відхилення впливної величини в межах робочих умов.\\[2mm]
Форма подання & Абсолютна, відносна або приведена.\\\bottomrule
\end{tabular}
\medskip

\normalsize
Приклад: передбачуваний температурний зсув нуля за сталої підвищеної температури й усталеного входу може одночасно бути систематичним, статичним та додатковим.
\medskip

Робочий температурний інтервал не означає збереження основної межі в усьому інтервалі без урахування додаткового впливу.
\end{frame}

\begin{frame}{Визначення: навчальна температурна модель}
Нормальний інтервал $[T_{\min};T_{\max}]$ лежить усередині робочого $[T_L;T_H]$. Для $T\in[T_L;T_H]$:
\[
d_T=\max(T_{\min}-T,\;0,\;T-T_{\max}),
\qquad B_T=k_Td_T,\quad B_{\mathrm{stat}}=B_0+B_T.
\]
$k_T\ge0$ має одиницю величини на градус, а $B_0$ — основна межа.
\medskip

\textbf{Приклад:} нормальні $15\ldots25\,{}^\circ\mathrm C$, робочі $-10\ldots60\,{}^\circ\mathrm C$, $B_0=0{,}10\unit{кПа}$, $k_T=0{,}01\unit{кПа}/{}^\circ\mathrm C$.
\[
T=40\,{}^\circ\mathrm C:\quad d_T=15\,{}^\circ\mathrm C,\quad B_T=0{,}15\unit{кПа},\quad B_{\mathrm{stat}}=0{,}25\unit{кПа}.
\]
Це навчальна модель, не апроксимація WIKA. Поза робочими умовами числову межу за цією формулою не заявляємо.
\end{frame}

\begin{frame}{Реальний паспорт: температура та виконання}
У WIKA S-20, паспорт PE 81.61, редакція 02/2026:
\begin{itemize}
\item температурний вплив описано окремо;
\item базова температура налаштування становить $15\ldots25\,{}^\circ\mathrm C$, є інші замовні варіанти;
\item допустимі параметри залежать від виконання.
\end{itemize}
Отже, дані конкретної модифікації та фактичні умови потрібні до розрахунку. Навчальний лінійний коефіцієнт із попереднього слайда не є числом із цього паспорта.
\medskip

Якщо комплексна характеристика вже охоплює певний температурний ефект, повторне додавання того самого внеску потребує усунення.
\src{https://www.wika.com/media/Data-sheets/Pressure/Pressure-sensors/ds_pe8161_en_co.pdf}{WIKA S-20, PE 81.61, 02/2026, с. 2–3}
\end{frame}

\begin{frame}{Визначення: динамічна похибка першого порядку}
Навчальна модель з одиничним статичним коефіцієнтом:
\[
\tau\frac{dy}{dt}+y=x,\qquad \tau>0,\qquad e=y-x.
\]
Оскільки $y=e+x$, отримуємо
\[
\tau\frac{de}{dt}+e=-\tau\frac{dx}{dt}.
\]
Для $x(t)=rt$, $x(0)=y(0)=0$:
\[
e(t)=-\tau r\left(1-e^{-t/\tau}\right),\qquad
\lim_{t\to\infty}e(t)=-\tau r.
\]
При $\tau=0{,}5\unit{с}$ і $r=2\unit{кПа/с}$ усталене відставання дорівнює $-1\unit{кПа}$. Частіше зчитування саме по собі не змінює $\tau$.
\end{frame}

\begin{frame}{Межа динамічного внеску та початковий стан}
Для тієї самої моделі:
\[
e(t)=e(0)e^{-t/\tau}-\int_0^t e^{-(t-s)/\tau}\dot x(s)\,ds.
\]
Якщо $|\dot x(s)|\le r_{\max}$, то
\[
|e(t)|\le |e(0)|e^{-t/\tau}+\tau r_{\max}(1-e^{-t/\tau}).
\]
За узгодженого початкового стану $e(0)=0$:
\[
|e(t)|\le\tau r_{\max}.
\]
Це властивість заданої моделі, а не універсальна межа будь-якого датчика. Паспортний час відповіді не можна автоматично ототожнювати з $\tau$ без визначеної процедури.
\end{frame}

\section{Нормування інструментальної похибки}
\begin{frame}{Визначення: навантаження джерела сигналу}
Модель: ненавантажена напруга $V_0$, послідовний опір $R_s\ge0$, вхідний опір приймача $R_{\mathrm{in}}>0$.
\[
V_{\mathrm{in}}=V_0\frac{R_{\mathrm{in}}}{R_s+R_{\mathrm{in}}},\qquad
e=V_{\mathrm{in}}-V_0=-V_0\frac{R_s}{R_s+R_{\mathrm{in}}}.
\]
Для $V_0\ne0$:
\[
\frac e{V_0}=-\frac{R_s}{R_s+R_{\mathrm{in}}}.
\]
\textbf{Навчальний приклад:} $V_0=5\unit{В}$, $R_s=100\unit{Ом}$, $R_{\mathrm{in}}=10\unit{кОм}$:
\[
V_{\mathrm{in}}\approx4{,}9505\unit{В},\qquad e\approx-49{,}5\unit{мВ}.
\]
Ці опори не є заявленою внутрішньою моделлю конкретного передавача.
\end{frame}

\begin{frame}{Перевірка навантаження за реальним паспортом}
Для виходу напруги WIKA S-20 паспорт задає:
\[
R_L>\frac{\text{максимальний вихідний сигнал}}{1\unit{мА}}.
\]
Для виконання $0\ldots10\unit{В}$:
\[
R_L>10\unit{кОм}.
\]
\begin{itemize}
\item Це умова застосування паспортного опису, а не значення внутрішнього вихідного опору.
\item Поза нею потрібна інша обґрунтована модель або зміна підключення.
\item Перед окремим додаванням навантажувального внеску перевіряємо, чи вже охоплює його комплексна межа.
\end{itemize}
\src{https://www.wika.com/media/Data-sheets/Pressure/Pressure-sensors/ds_pe8161_en_co.pdf}{WIKA S-20, PE 81.61, 02/2026, с. 7: Load in Ω}
\end{frame}

\begin{frame}{Визначення: неінформативна пульсація}
Навчальна модель: тиск кодує постійний рівень $Sp$, а пульсація його не кодує:
\[
V(t)=Sp+A\sin(\omega t+\varphi),\qquad S\ne0.
\]
При миттєвому зчитуванні в моменти $t_i$:
\[
\widehat p_i=\frac{V(t_i)}S,\qquad
e_{p,i}=\frac A S\sin(\omega t_i+\varphi),\qquad |e_{p,i}|\le\frac{|A|}{|S|}.
\]
\begin{itemize}
\item Однакова фаза зчитування може зберігати зміщення навіть при багатьох повторах.
\item Випадкову фазу та відповідний розподіл потрібно обґрунтувати окремо.
\item Інший приклад: синфазна напруга на різницевому вході; перевіряють її допустиму область та вплив приймача.
\end{itemize}
Напруга живлення є впливною величиною; породжена нею пульсація вже належить властивостям вихідного сигналу.
\end{frame}

\begin{frame}{Навчальний бюджет: два окремі підсумки}
\small
\begin{center}
\begin{tabular}{lrrrrr}
\toprule
Внесок & Основний & Темпер. & Динаміч. & Навантаж. & Пульсація\\\midrule
мВ & 10 & 6 & 4 & 5 & 2\\\bottomrule
\end{tabular}
\end{center}
\normalsize
Припускаємо різні залишкові ефекти без подвійного врахування:
\[
B_\Sigma=10+6+4+5+2=27\unit{мВ}.
\]
Окремо задано незалежні випадкові внески $\sigma_1=3$, $\sigma_2=4\unit{мВ}$:
\[
\sigma_\Sigma=\sqrt{3^2+4^2}=5\unit{мВ}.
\]
Якщо випадкова сума нормальна і $\sigma_\Sigma$ відома:
\[
L=B_\Sigma+1{,}96\sigma_\Sigma=36{,}8\unit{мВ}
\quad\Longleftrightarrow\quad L_p=0{,}368\unit{кПа}
\]
для каналу $S=0{,}1\unit{В/кПа}$. Модельне покриття близько 95\% або більше, але не безумовна гарантія.
\end{frame}

\begin{frame}{Структура перевірного бюджету}
\begin{itemize}
\item Кожний рядок містить ефект, категорію, тип характеристики, значення, одиницю, умови та джерело.
\item Знакові оцінки, межі та СКВ зберігаємо окремо.
\item Перевіряємо, що комплексна основна характеристика не охоплює вже доданий окремий ефект.
\item Якщо залежність випадкових внесків невідома, RSS із нульовими коваріаціями не заявляємо як підтверджену оцінку.
\item Поза робочими умовами або за відсутності потрібного параметра повертаємо статус неповного обґрунтування.
\end{itemize}
Повнота бюджету залежить від моделі та перевірених умов, а не від кількості рядків або величини підсумку.
\end{frame}

\begin{frame}{Офіційні питання: блоки 1–2}
\small
\begin{enumerate}
\setcounter{enumi}{25}
\item Постановка задачі підсумовування похибок вимірювань.
\item Визначення сумарної систематичної, випадкової та повної похибки вимірювань.
\item Нормування характеристик похибок вимірювань.
\item Визначення, суть і мета нормування МХ ЗВТ.
\item Нормовані метрологічні характеристики ЗВТ призначені для визначення результату вимірювання.
\end{enumerate}
\medskip

Блок 1 покриває 26–28, блок 2 — 29–30. Формулювання звірено з офіційним переліком у DOCX 2020 року.
\end{frame}

\begin{frame}{Офіційні питання: блоки 3–4}
\small
\begin{enumerate}
\setcounter{enumi}{30}
\item Класифікаційні ознаки похибок ЗВТ.
\item Форми відображення кількісних характеристик похибок ЗВТ.
\item Систематичні, випадкові, статичні, динамічні похибки ЗВТ.
\item Основна і додаткова похибки ЗВТ.
\item Вибір комплексу нормованих складових інструментальної похибки вимірювань.
\item Нормування основної, додаткової, динамічної похибки ЗВТ.
\item Нормування похибки взаємодії та не інформованих параметрів вихідних сигналів ЗВТ.
\end{enumerate}
\medskip

Блок 3 покриває 31–34, блок 4 — 35–37. У питанні 37 збережено текст оригіналу; змістовний термін у лекції — «неінформативні параметри».
\end{frame}

\begin{frame}{Джерела та подальша робота}
\small
\begin{itemize}
\item \href{https://www.nist.gov/pml/nist-technical-note-1297/nist-tn-1297-appendix-law-propagation-uncertainty}{NIST TN 1297, Appendix A: поширення невизначеності}.
\item \href{https://jcgm.bipm.org/vim/en/4.26.html}{JCGM VIM 4.26: maximum permissible measurement error}.
\item \href{https://www.fluke.com/en-us/product/electrical-testing/digital-multimeters/fluke-117}{Fluke 117: специфікації діапазонів DC V}.
\item \href{https://www.wika.com/media/Data-sheets/Pressure/Pressure-sensors/ds_pe8161_en_co.pdf}{WIKA S-20, PE 81.61, 02/2026: температурний вплив і навантаження}.
\item \href{https://www.ni.com/en/shop/data-acquisition/sensor-fundamentals/measuring-pressure-with-bridge-based-and-other-pressure-sensors.html}{NI: канал вимірювання тиску з містковим сенсором}.
\end{itemize}
Використано тематичний план курсу та наданий «Конспект МТВП»; припущення статистичних оцінок зазначено явно.
\medskip

Наступні чотири слайди містять розгорнуті Python-завдання без готових розв’язків.
\end{frame}

\appendix
\begin{frame}{Практичне завдання 1: підсумовування внесків}
\small
\textbf{Вхід, кПа:} межі \texttt{[0.12, 0.08, 0.05]}; СКВ незалежних внесків \texttt{[0.03, 0.04, 0.02]}; знакові оцінки \texttt{[-0.06, 0.02]}.
\begin{enumerate}
\item Створіть \texttt{bound\_sum(bounds)} для граничної суми та \texttt{rss(sigmas)} для СКВ суми незалежних внесків.
\item Окремо обчисліть алгебраїчну суму знакових оцінок і підпишіть усі три результати.
\item Для меж і СКВ вимагайте невід’ємних скінченних чисел; знакові оцінки можуть бути від’ємними.
\item Порожній список відхиляйте як відсутність бюджету; початкові списки не змінюйте.
\end{enumerate}
\textbf{Перевірки:} один внесок; нульовий внесок; перестановка елементів; від’ємна межа; нескінченність; порожній список.
\medskip

\textbf{Подати:} функції, три підписані підсумки й одне речення про те, чому RSS не є гарантованою межею. Достатньо \texttt{sum}, циклів та \texttt{math.sqrt}/\texttt{math.isfinite}.
\end{frame}

\begin{frame}{Практичне завдання 2: паспорт у словнику}
\small
\textbf{Навчальні дані:} режим DC; діапазон $[0;10]\unit{В}$; $a=0{,}002$, $d=3$, $q=0{,}001\unit{В}$; покази \texttt{[0.2, 5.0, 9.8]} В.
\begin{enumerate}
\item Створіть словник параметрів та функцію оцінки $B(x)=a|x|+dq$ для заданого показу.
\item Перевіряйте наявність потрібних полів, режим, скінченність, діапазон, $a,d\ge0$ і $q>0$.
\item Повертайте значення межі разом зі статусом; поза діапазоном не видавайте формулу як чинну специфікацію.
\item Інші паспортні умови у вправі приймаються виконаними; явно зазначте це у звіті.
\end{enumerate}
\textbf{Перевірки:} нуль; верхня межа; значення вище діапазону; відсутнє поле; неправильний режим. Якщо додаєте $B/|x|$, для $x=0$ повертайте \texttt{None}.
\medskip

\textbf{Подати:} словник, функцію й таблицю результатів з одиницями. Ці коефіцієнти не належать реальному Fluke 117.
\end{frame}

\begin{frame}{Практичне завдання 3: температурні умови}
\small
\textbf{Вхід:} $B_0=0{,}20\unit{кПа}$, $k_T=0{,}015\unit{кПа}/{}^\circ\mathrm C$; нормальні $[18;28]$, робочі $[-10;60]\,{}^\circ\mathrm C$; \texttt{T = [20, 28, 35, 60, 65]}.
\begin{enumerate}
\item Функція приймає параметри та температуру; перевіряє порядок меж і вкладеність нормального інтервалу в робочий.
\item Усередині робочих умов обчисліть відстань до нормального інтервалу $d_T$, додаткову межу $B_T=k_Td_T$ та суму $B_0+B_T$.
\item Поза робочим інтервалом повертайте \texttt{None} для числової оцінки й пояснювальний статус; межі інтервалів включені.
\item Виведіть для кожної температури $d_T$, $B_T$, суму та статус; підпишіть одиниці.
\end{enumerate}
\textbf{Перевірки:} обидві нормальні межі; обидві робочі межі; температура всередині нормального інтервалу; некоректний порядок меж; температура поза робочим інтервалом.
\medskip

\textbf{Подати:} функцію, таблицю та пояснення відмінності між працездатністю і збереженням основної межі.
\end{frame}

\begin{frame}{Практичне завдання 4: бюджет із категоріями}
\small
\textbf{Вхід, мВ:} межі: основна 12, температура 8, динаміка 6, навантаження 4; СКВ: 2 та 3; $k=1{,}96$.
\begin{enumerate}
\item Подайте внески списком словників: \texttt{id}, категорія, тип \texttt{bound}/\texttt{sigma}, значення, одиниця; додайте перелік уже охоплених основною межею ефектів.
\item Перевірте одиниці, скінченність, невід’ємність, повтори ідентифікаторів та подвійне врахування ефектів.
\item Обчисліть $B_\Sigma$; RSS повертайте лише за явно прийнятої незалежності, інакше \texttt{None} та пояснення.
\item $L=B_\Sigma+k\sigma_\Sigma$ обчислюйте лише за також прийнятих нормальності суми й відомих СКВ; підпишіть модельний, а не гарантований зміст.
\end{enumerate}
\textbf{Перевірки:} різні одиниці; дубль ефекту; ефект уже охоплено основною межею; невідома залежність; відсутня категорія.
\medskip

\textbf{Подати:} функцію, підсумки та звіт умов; початкові дані збережіть, використайте функції задачі 1.
\end{frame}
\end{document}
